\section{Visible repetition and genealogical differentiation within a region}
\label{md:ejemplo:region}
The biographies in \eqref{act21:eq:biografias-tpk-tres-caracteres}
make it possible to specify the difference between a repetition of digits and a
return of the state. For the channel recognized as \(e\), the
calendar \eqref{act21:eq:calendario-0011-interno} produces
\[
 201101\mid121221\mid102011\mid012222\mid102011.
\]
Lifting precedes the decimal readout. The same matrices act
on the other two channels, preserving their initial emissions and their
distinct regional states. This common structure makes it possible to compare
the biographies without identifying their coordinates.

Let \(w\) be a concatenation of \(L\) trits, and let \(N\) be its integer
value in base three. The corresponding cylinder is
\[
 I_w=[N/3^L,(N+1)/3^L).
\]
For \(d\) decimal places, the prefix is stable throughout the
cylinder exactly when
\[
 \left\lfloor\frac{10^dN}{3^L}\right\rfloor
 =\left\lfloor\frac{10^d(N+1)-1}{3^L}\right\rfloor.
\]
This applies Lemma~\ref{act21:censo:estabilidad} to the decimal
readout. The half-open right endpoint explains the term \(-1\)
in the integer division. Adding the integer part of the regional chart,
the five successive readouts are
\begin{center}
\begin{tabular}{@{}rl@{}}
\toprule
Tritic length & Stable prefix of the regional readout\\
\midrule
6 & \(2.71\ldots\)\\
12 & \(2.71828\ldots\)\\
18 & \(2.7182818\ldots\)\\
24 & \(2.7182818284\ldots\)\\
30 & \(2.71828182845904\ldots\)\\
\bottomrule
\end{tabular}
\end{center}
At length twenty-four, one has
\(N=202864003874\) and \(3^{24}=282429536481\).
At length thirty, one obtains
\(N=147887858824447\) and \(3^{30}=205891132094649\).
The preceding integer divisions certify the prefixes in the table.
In particular, the first twelve fractional digits group as
\[
 718281828459
 =7\,\underbrace{1828\,1828}_{\text{positional repetition}}\,459.
\]
The twenty-four-trit cylinder already fixes both copies and the
following digit. This localization distinguishes the decimal episode from the
reappearance of the tritic block \(102011\), located in the
third and fifth positions: their cuts and their readout bases differ.

Repetition of the visible block does not identify its prestates either.
The prestates modulo twenty-seven retained in the development of
this biography are
\[
 (25,11,7,17,8,16),\qquad (7,8,4,23,26,7),
\]
with respective subsequent quotients
\[
 (6,13,9,2,13,1),\qquad (15,0,3,19,23,25).
\]
Equality of the emission thus coexists with a differentiated genealogy.
Its explanatory role is specific: observing a repeated
block retains less information than the enriched continuation
that produces it. To attribute a future difference exclusively
to memory, the phases and comparison charts must also
be retained.

This episode relates three levels: the transitions produce the
tritic chain, the cylinders determine the stable digits, and the
state retains the provenance of a repeated emission. Arbitrary
depth belongs to the prolongation proved
earlier; this finite readout specifies one episode of that
biography. The regularity of interest lies in the
compatibility between levels, not in prolonging a decimal fragment
indefinitely: a third copy of \(1828\) would begin with
\(1\), whereas the continuation already produced begins with \(4\).
